% ---
% id: protocol-compiler/req:witness-refusal-meaning
% kind: req
% title: "Reporting a refusal for want of a feasibility witness"
% summary: >-
%   A refusal at the feasibility check is never reported to a user as a defect
%   of the protocol: it means only that no feasible state was shown, and a
%   well-formed declaration can be refused there.
% prominence: standard
% status: accepted
% uses: [rem:completeness, thm:validator-soundness]
% assumes: []
% refs: []
% provenance: newly-derived
% lean:
% ---

Check V7 is the one check that can refuse a \tref{well-formed declaration}
(Theorem~\ref{thm:validator-soundness}\,(iii)): its refusal means that no feasible state was
shown, not that none exists (Remark~\ref{rem:completeness}).

\begin{requirement}\label{req:witness-refusal-meaning}
The engine never reports a refusal with \iface{no_feasibility_witness} to a
user as a defect in the protocol.
\end{requirement}
